\chapter{Introdução}
%---------------------------------------------------------------------------------------

As eleições municipais brasileiras ocorrem a cada quatro anos, e os prefeitos que as disputam, diretamente ou por meio de aliados, controlam orçamentos responsáveis por serviços de grande visibilidade para a população, como saúde, educação, obras urbanas e transporte. A literatura sobre ciclos políticos orçamentários investiga se, e de que forma, os governantes usam esse controle para influenciar o resultado das urnas. Na abordagem oportunista inaugurada por \textcite{nordhaus1975}, o governante expande a atividade econômica antes da eleição para obter votos; nos modelos de sinalização de \textcite{rogoff1990}, com eleitores racionais, o ciclo surge porque o governante usa gastos de rápida percepção para sinalizar competência. \textcite{drazen2010} mostram que o ciclo pode ocorrer mesmo sem a geração de déficits, pela recomposição do orçamento em favor das despesas mais visíveis ao eleitor.

No Brasil, a evidência empírica para os municípios é ampla. \textcite{sakurai2011} encontram aumento das despesas totais e correntes nos anos eleitorais, e \textcite{sakurai2009}, ao analisar as despesas por função entre 1990 e 2005, mostra que o ano eleitoral eleva o gasto em áreas como saúde e saneamento, habitação e urbanismo e transporte. A literatura também indica que o alinhamento político importa: \textcite{ferreira2007} mostram que municípios cujo prefeito pertence à coligação que elegeu o governador ou ao mesmo partido do presidente recebem mais transferências voluntárias. Desde o período analisado por esses estudos, porém, o contexto mudou: os municípios passaram a informar suas contas em um sistema padronizado, o Sistema de Informações Contábeis e Fiscais do Setor Público Brasileiro (Siconfi), e o período recente reuniu uma recessão, uma pandemia que coincidiu com a eleição de 2020 e mudanças bruscas na coalizão do governo federal, que alteraram a proporção de prefeitos alinhados ao presidente.

Diante desse contexto, o problema desta pesquisa pode ser formulado na seguinte pergunta: de que forma o calendário eleitoral e o alinhamento político do prefeito com o governador e com o presidente afetam o volume e a composição das despesas dos municípios brasileiros entre 2013 e 2025? Com base na literatura, formulam-se duas hipóteses. A primeira é que, no ano eleitoral, os prefeitos recompõem o orçamento em favor das despesas mais visíveis ao eleitor, como obras urbanas, transporte e saúde, em detrimento das menos visíveis, como as administrativas. A segunda é que o alinhamento do prefeito com o governador e com o presidente está associado a uma despesa municipal maior, em razão do maior acesso a transferências voluntárias. O objetivo geral é estimar como as despesas municipais respondem ao calendário eleitoral e ao alinhamento político do prefeito no período de 2013 a 2025. Os objetivos específicos são: (i) estimar o efeito do ano eleitoral e do ano pré-eleitoral sobre a despesa total e sobre oito grupos de funções de governo; (ii) verificar se esse efeito se repete nas eleições de 2016 e de 2024; (iii) estimar o efeito do alinhamento do prefeito com o governador e com o presidente sobre as mesmas despesas; e (iv) comparar os resultados com os de \textcite{sakurai2009}.

Para isso, foi montado um painel de cerca de 5,6 mil municípios com dados fiscais do Siconfi, eleitorais do Tribunal Superior Eleitoral (TSE) e demográficos e econômicos do Instituto Brasileiro de Geografia e Estatística (IBGE). As despesas foram agrupadas por função, seguindo \textcite{sakurai2009}, e analisadas em dois modelos de efeitos fixos: o primeiro mede o ciclo eleitoral, com erros-padrão robustos a choques comuns a todos os municípios; o segundo mede o efeito do alinhamento político, comparando, dentro do mesmo ano, prefeitos alinhados e não alinhados. A despesa entra em logaritmo, de modo que os efeitos se leem como variações percentuais, e a robustez dos resultados é avaliada em dezoito testes, que variam a forma funcional, a amostra, os controles, o tratamento dos dados e o tipo de erro-padrão.

Os resultados indicam que o ciclo eleitoral aparece sobretudo na composição do gasto: no ano eleitoral, as despesas com Habitação e Urbanismo, Transporte e Saúde e Saneamento aumentam nas duas eleições analisadas, enquanto Administração e Educação e Cultura perdem participação na despesa total. A despesa total é cerca de 5\% maior no ano eleitoral, mas esse efeito é mais forte em 2024; o alinhamento com o governador está associado a uma despesa total cerca de 0,4\% maior, e o alinhamento com o presidente não tem efeito sobre ela.

O tema é relevante porque os municípios são responsáveis pela oferta direta de serviços como saúde, educação e infraestrutura urbana, de modo que distorções eleitorais na alocação desses recursos afetam diretamente a população e podem comprometer o planejamento de longo prazo. Além disso, o trabalho contribui para a literatura ao atualizar a evidência sobre ciclos políticos orçamentários nos municípios brasileiros para um período recente, com praticamente todos os municípios do país e a classificação funcional atual. O trabalho também separa o efeito do calendário eleitoral do efeito do alinhamento político em modelos distintos, o que evita confundir o alinhamento com o presidente com as diferenças entre governos, e avalia a sensibilidade dos resultados a um conjunto amplo de decisões empíricas.

Além desta introdução, o trabalho tem outros cinco capítulos. O capítulo~\ref{cap:modelos} apresenta os modelos teóricos dos ciclos políticos e a evidência empírica, e o capítulo~\ref{cap:mecanismos} discute mecanismos específicos do caso brasileiro, como o federalismo fiscal, o uso de restos a pagar e as regras eleitorais. O capítulo~\ref{cap:metodologia} descreve os dados e o método, o capítulo~\ref{cap:resultados} apresenta os resultados, e o capítulo~\ref{cap:conclusao} traz a conclusão. O \autoref{ap:robustez} reúne os testes de robustez.
